\chapter{Literature Review} \label{chap:sota}

\section*{}
==============

Neste capítulo é descrito o estado da arte e são
apresentados trabalhos relacionados para mostrar o que existe no
mesmo domínio e quais os problemas em aberto.
Deve deixar claro que existe uma oportunidade de desenvolvimento que
cobre alguma falha concreta.

O capítulo deve também efetuar uma revisão tecnológica às principais
ferramentas utilizáveis no âmbito do projeto, justificando futuras
escolhas.

==============

\section{Introdução}

Music Information Retrieval (MIR) is the field of research that focuses on the
development of algorithms and systems to analyze, retrieve, and organize music data.
MIR bridges a wide range of disciplines, including: 
musicology, signal processing, machine learning, and human-computer interaction. 


==============
CONTINUAR 

@mullerFundamentalsMusicProcessing2015
==============

\section{Visualization Types}\label{sec:dialecto}


We can sum up all music related visualizations into three categories:

    1- Sheet Music and Score Images
    2- Symbolic Representations
    3- Audio Representations

\subsection{Sheet Music and Score Images}

Sheet music describes musical work using a formal language based on musical symbols and letters,
thus, this kind of music representations always require  a level of music literacy from the 
performing musician. Furthermore sheet music usally does not induce explicit information on how 
a given piece should be performed, it is intened for the musician to have a certain 
level of interpretation liberaty. This can include the variance of tempo, articulation and 
dynamics amoung others.

Sheet music also comes in a variety of forms that vary throughout History and cultures. 
For our porpuses we are going to focus on the western standard of music notation. And still, within
wester music, there can be vastly differnt forms of notation, however, most are based on a staff
which sets five horizontal lines and four spaces, each representing a different musical pitch.
Musical relevant symbols are put inside the staff to denote pitch, temporal and expressive indications. 

@mullerMusicRepresentations2015a
\subsection{Symbolic Representations}

Symbolic Representations are machine-readable formats where 
musical events and expressions are explicitly encoded.

These come in various formats, such as MIDI and MusicXML.

Visualy we can identify two main applications for this category, 
these include: Piano-Roll representations and digital score representations. 

Historically piano-roll representations come from so-called \textbf{player pianos},
these were self playing  pianos that became quite popular in the late 19th to early 20th century. 
This was achieved by a pneumatic mechanism that would hit keys and pedal of the piano,
through a continous roll of paper with perforations. The roll would move over a reading system 
(known as a tracker bar) that would  trigger the musical actions accordingly. 

This concept of a continuos and explicit representation of music performance would later prove usefull
for digital applications. Such is the case for the Piano-Roll.

The \textbf{Piano-Roll} is a stable on any DAW as a visual and interactive tool to write and edit musical events
in the form of MIDI notes and atributes (such as velocity and other expressive controls).

\textbf{MusicXML} refers to the XML-based language to transcribe score music into the digital domain. 
There are multiple sub-varients of the MusicXML format for specific aplications, 
however all share an XML based langue underneath. 

MusicXML was developed to serve as a universal format for storing and sharing music scores
across differnt music notation applications. Music notation digital editors then 
employ proprietary variants of this format for addressing the program's specific charecteristics.
In some cases these may not even be xml based, but MusicXML format can be considered the universal format.

These include:
    .mscz and .mscx - These are Musescore specific formats named MuseScore format and Uncompressed MuseScore format respectively. Both xml based.
    .dora - "Dorico" specific format (non xml based)
    .sib  - "Sibelius" proprietary format  of company Avid (also non xml based)
    .mei  - The "Music Encoding Initiative" developed format for accademic and schollar research. XML based.

\subsection{Audio Representations}
Audio representations are graphs and/or plots that intend to represent sound. 
These may not be for sctrictly representing music, the intent is to translat sound
from the audible to the visual domain. 
Audio is a term that refers to an acoustic sound that is turned into an electrical signal. Most audio
representations, are thus digitaly generated through applying mathematical equations
to the digitalized audio signal, such as the Fourier Transform. 

We can infer at least 3 components in audio. These are, time, frequency and amplitude.
ways of visualy represening these elements, not all audio representations try to describe the three. 

Waveform - Is a the most comum and one of the simplest ways to represent audio.
It consists of a single line on a graph that represents a given sound source, this 2d graph
has amplitude as it's y-axis and time as it's x-axis. This representation does not provide data into frequency
it only provides information for amplitude (how loud a sound is) and respective time data.

For displaying the three sound components (time, amplitude and frequency/pitch) the graph needs
to some how represnt three distict dimensions.  This can be done trhough developing a 3D graph 
with axis for Amplitude, Time and Frequency. 

The spectrogram achieves this by add by having on the y axis the frequency, time on the x-axis 
and color as substitute for the third axis (amplitude), and in this way displaying 
all three sound physical parameters in a single 2D visual plot. 

@henriqueAcusticaMusical
@jungUnifiedCrossmodalTranslation2026a; @mullerMusicRepresentations2015a
\section{MIR Visualization Tools}\label{sec:dialecto}

\section{Annotation and Evaluation Tasks in MIR}\label{sec:dialecto}


\section{Resumo ou Conclusões}
